\subsection{Функции}

\textbf{Мотивирующий вопрос:} Как не писать один и тот же код для взлета в разных частях программы?

Функция — это именованный блок кода, который можно вызывать многократно. Функции делают код чище и понятнее.

Ключевое слово \texttt{def} объявляет функцию. \texttt{return} возвращает результат.

\begin{codefile}[python]{examples/python/07_functions.py}
\end{codefile}

\begin{infobox}[Область видимости]
Переменные, созданные внутри функции, не видны снаружи. Это называется локальной областью видимости.
\end{infobox}

\begin{taskbox}[Конвертер]
Напишите функцию \texttt{celsius\_to\_fahrenheit(c)}, которая переводит градусы Цельсия в Фаренгейты по формуле $F = C \times 1.8 + 32$.
\end{taskbox}
